% between spacetime representation theory and local gauge symmetry.

\section{作用量、场方程与 Noether 守恒律}\label{sec:action-noether-dp61}

表示论告诉我们场可以怎样变换，却还没有规定哪一种运动由自然界实现。作用量把“允许的场表示”与“实际的动力学”连接起来，而 Noether 定理进一步把连续对称性转化为守恒量，为随后把整体相位对称局域化做好准备。

\subsection{从局域拉氏量到场方程}

质点力学里，Euler--Lagrange 方程来自“真实轨迹使作用量驻值”。场论没有改变这个原则，只是把有限个坐标 \(q_a(t)\) 换成每个时空点都有取值的场 \(\Phi_A(x)\)。因此这里首先要回答的不是“场方程长什么样”，而是怎样把普通变分法推广到局域场，并看清边界项为何决定运动方程与守恒流。


令 $\Phi_A(x)$ 表示一组经典场，$A$ 只用于枚举不同场分量。全文取 $x^0=ct$，因此 $\dd^4 x=c\,\dd t\,\dd^3 r$。若拉氏量密度只依赖同一点的场及其一阶导数，作用量定义为
\begin{equation}
 \boxed{
 S[\Phi]=\frac1c\int_{\Omega}\dd^4 x\,
 \Lag\!\left(\Phi_A,\partial_\mu\Phi_A\right).}
 \label{eq:field-action-general-dp61}
\end{equation}
这里 $\Omega$ 是所考察的时空区域；$[\Lag]=\mathrm{J\,m^{-3}}$，所以 $[S]=\mathrm{J\,s}$。作用量 $S$ 是把整段时空历史映射到一个具有作用量量纲标量的泛函；其驻值条件决定经典运动方程。

对场作任意小变分 $\Phi_A\to\Phi_A+\delta\Phi_A$。只要在边界 $\partial\Omega$ 上固定场，即 $\delta\Phi_A|_{\partial\Omega}=0$，驻值条件 $\delta S=0$ 就要求
\begin{equation}
 \boxed{
 \partial_\mu\frac{\partial\Lag}{\partial(\partial_\mu\Phi_A)}
 -\frac{\partial\Lag}{\partial\Phi_A}=0.}
 \label{eq:field-euler-lagrange-dp61}
\end{equation}
这是场的 Euler--Lagrange 方程。它与质点力学中的 Euler--Lagrange 方程完全同源；差别只在于离散自由度指标被时空点和场分量共同替代。



\subsection{连续对称怎样变成守恒流}

场方程告诉我们给定初态以后怎样演化，但还没有解释为什么某些组合在整个演化中保持不变。连续对称性的作用正是把“变换后物理不变”转化成一个局域连续性方程。与只记住“对称性对应守恒量”相比，Noether 定理更强：它直接给出守恒流的表达式，并说明守恒荷只能通过边界流动而不能在体内凭空产生。

先考虑不移动时空坐标的连续内部变换。令 $\epsilon$ 是一个与 $x$ 无关的无穷小常数，并写
\begin{equation*}
 \delta\Phi_A=\epsilon\,\Delta_A(\Phi).
\end{equation*}
若拉氏量的变化最多是一个四维散度，
\begin{equation}
 \delta\Lag=\epsilon\,\partial_\mu K^\mu,
 \label{eq:noether-total-derivative-dp61}
\end{equation}
则作用量在固定边界条件下不变。利用 Euler--Lagrange 方程把体项消去以后，剩余变化必须写成局域连续性方程
\begin{equation}
 \boxed{
 \partial_\mu j^\mu=0,
 \qquad
 j^\mu
 =\sum_A
 \frac{\partial\Lag}{\partial(\partial_\mu\Phi_A)}\,\Delta_A
 -K^\mu.}
 \label{eq:noether-current-general-dp61}
\end{equation}
这就是固定坐标内部连续对称性的 Noether 定理。\eqnrefs{eq:noether-current-general-dp61} 比“对称性对应守恒量”更具体：它告诉我们守恒流怎样从拉氏量逐项计算出来。

若把该流写成 $j^\mu=(c\rho_j,\mathbf j)$，则
\begin{equation*}
 \partial_\mu j^\mu=0
 \quad\Longleftrightarrow\quad
 \partial_t\rho_j+\boldsymbol\nabla\!\cdot\mathbf j=0.
\end{equation*}
对任意固定空间体积 $V$ 积分并用 Gauss 定理，得到
\begin{equation}
 \boxed{
 \frac{\dd}{\dd t}\int_V\dd^3 r\,\rho_j
 =-\oint_{\partial V}\mathbf j\cdot\dd\mathbf S.}
 \label{eq:noether-flux-balance-dp61}
\end{equation}
左边是区域内守恒荷的变化率，右边是穿过边界的净流出。若取全空间且场在无穷远处衰减得足够快，边界通量消失，于是总荷
\begin{equation*}
 Q_j=\int\dd^3 r\,\rho_j
\end{equation*}
不随时间变化。局域连续性方程因此比“总量守恒”更强：它还规定守恒量只能通过边界输运，不能在体内凭空产生或消失。

\begin{exbox}[自由 Dirac 场的整体 $U(1)$]
自由 Dirac 拉氏量为
\begin{equation*}
 \Lag_D
 =\bar\psi\left(\ii\hbar c\gamma^\mu\partial_\mu-\me c^2\right)\psi.
\end{equation*}
它在常数相位变换 $\psi\to \ee^{\ii\theta}\psi$、$\bar\psi\to\bar\psi \ee^{-\ii\theta}$ 下不变。Noether 流的整体倍数取决于生成元的归一化；选择粒子数流的常用归一化后，
\begin{equation}
 \boxed{n^\mu=c\,\bar\psi\gamma^\mu\psi,
 \qquad \partial_\mu n^\mu=0.}
 \label{eq:dirac-number-current-dp61}
\end{equation}
于是
\begin{equation*}
 n^0=c\psi^\dagger\psi,
 \qquad
 N=\int\dd^3 r\,\psi^\dagger\psi.
\end{equation*}
若该 Dirac 场携带电荷 $\qe$，物理电流就是
\begin{equation}
 \boxed{j^\mu=\qe n^\mu=\qe c\,\bar\psi\gamma^\mu\psi.}
 \label{eq:dirac-electric-current-from-noether-dp61}
\end{equation}
整体 $U(1)$ 对称取 $\theta$ 为时空常数，因此普通导数不产生额外项。令相位依赖时空位置 $\theta=\theta(x)$ 后，$\partial_\mu\psi$ 会额外产生与 $\partial_\mu\theta$ 成正比的项。抵消该项需要引入电磁四势，并用协变导数取代普通导数。对应关系为
\[
 \text{整体相位对称性}
 \longrightarrow \text{Noether 守恒流}
 \longrightarrow \text{局域化相位}
 \longrightarrow \text{规范联络 }A_\mu.
\]


\medskip

一般 Noether 流~\eqnrefs{eq:noether-current-general-dp61} 已经建立。对自由 Dirac 拉氏量，只需把整体相位变换 $\delta\psi=\ii\epsilon\psi$、$\delta\bar\psi=-\ii\epsilon\bar\psi$ 代入该公式。质量项与动能项的相位变化彼此抵消，所以 $K^\mu=0$；再按“每个粒子携带单位数荷”的约定重标度生成元，就得到\eqnrefs{eq:dirac-number-current-dp61}。

取无穷小整体变换
\begin{equation*}
 \delta\psi=\ii\epsilon\psi,
 \qquad
 \delta\bar\psi=-\ii\epsilon\bar\psi,
 \qquad \partial_\mu\epsilon=0.
\end{equation*}
质量项的变化为
\begin{equation*}
 \delta(\bar\psi\psi)
 =(-\ii\epsilon\bar\psi)\psi
 +\bar\psi(\ii\epsilon\psi)=0.
\end{equation*}
动能项中
\begin{align*}
 \delta\!\left(\bar\psi\gamma^\mu\partial_\mu\psi\right)
 &=(-\ii\epsilon\bar\psi)\gamma^\mu\partial_\mu\psi
 +\bar\psi\gamma^\mu\partial_\mu(\ii\epsilon\psi),\\
 &=-\ii\epsilon\bar\psi\gamma^\mu\partial_\mu\psi
 +\ii\epsilon\bar\psi\gamma^\mu\partial_\mu\psi=0,
\end{align*}
所以 $K^\mu=0$。又因为
\begin{equation*}
 \frac{\partial\Lag_D}{\partial(\partial_\mu\psi)}
 =\ii\hbar c\,\bar\psi\gamma^\mu,
 \qquad
 \frac{\partial\Lag_D}{\partial(\partial_\mu\bar\psi)}=0,
\end{equation*}
Noether 公式给出的原始流为
\begin{equation*}
 j_{\rm raw}^\mu
 =(\ii\hbar c\bar\psi\gamma^\mu)(\ii\psi)
 =-\hbar c\bar\psi\gamma^\mu\psi.
\end{equation*}
连续群生成元可以整体重标度；相应 Noether 流也按同一常数重标度。把生成元归一化成“每个粒子携带单位数荷”，定义
\begin{equation*}
 n^\mu\equiv-\frac1\hbar j_{\rm raw}^\mu
 =c\bar\psi\gamma^\mu\psi.
\end{equation*}
于是 $\partial_\mu n^\mu=0$。时间分量使用 $\bar\psi=\psi^\dagger\gamma^0$：
\begin{equation*}
 n^0=c\bar\psi\gamma^0\psi
 =c\psi^\dagger\psi.
\end{equation*}
故数荷为 $N=c^{-1}\int\dd^3 r\,n^0=\int\dd^3 r\,\psi^\dagger\psi$。
\end{exbox}

\begin{derivation}[从\eqnrefs{eq:field-action-general-dp61}到\eqnrefs{eq:field-euler-lagrange-dp61}]
正文\eqnrefs{eq:field-action-general-dp61} 把场的整段时空历史映射成作用量。对场做任意小变分时，唯一需要处理的是含 $\partial_\mu\delta\Phi_A$ 的项；把它分部积分后，边界条件 $\delta\Phi_A|_{\partial\Omega}=0$ 消去表面项，而区域内部的 $\delta\Phi_A(x)$ 又可以独立任取，因此体积分系数必须逐点为零。这正给出正文\eqnrefs{eq:field-euler-lagrange-dp61}。

从定义出发，
\begin{align*}
 \delta S
 &=\frac1c\int_\Omega\dd^4 x\,
 \left[
 \frac{\partial\Lag}{\partial\Phi_A}\delta\Phi_A
 +\frac{\partial\Lag}{\partial(\partial_\mu\Phi_A)}
 \delta(\partial_\mu\Phi_A)
 \right],\\
 &=\frac1c\int_\Omega\dd^4 x\,
 \left[
 \frac{\partial\Lag}{\partial\Phi_A}\delta\Phi_A
 +\frac{\partial\Lag}{\partial(\partial_\mu\Phi_A)}
 \partial_\mu\delta\Phi_A
 \right].
\end{align*}
第二项使用乘积求导反解：
\begin{align*}
 \frac{\partial\Lag}{\partial(\partial_\mu\Phi_A)}
 \partial_\mu\delta\Phi_A
 &=\partial_\mu\!\left[
 \frac{\partial\Lag}{\partial(\partial_\mu\Phi_A)}
 \delta\Phi_A\right]
 -\partial_\mu\!\left[
 \frac{\partial\Lag}{\partial(\partial_\mu\Phi_A)}\right]
 \delta\Phi_A.
\end{align*}
代回以后，
\begin{align*}
 \delta S
 &=\frac1c\int_\Omega\dd^4 x\,
 \left[
 \frac{\partial\Lag}{\partial\Phi_A}
 -\partial_\mu\frac{\partial\Lag}{\partial(\partial_\mu\Phi_A)}
 \right]\delta\Phi_A\\
 &\quad+\frac1c\int_\Omega\dd^4 x\,
 \partial_\mu\!\left[
 \frac{\partial\Lag}{\partial(\partial_\mu\Phi_A)}
 \delta\Phi_A\right].
\end{align*}
四维 Gauss 定理把最后一项变成 $\partial\Omega$ 上的面积分；由于边界上 $\delta\Phi_A=0$，该项消失。于是
\begin{equation*}
 \delta S
 =\frac1c\int_\Omega\dd^4 x\,
 \mathfrak E_A(x)\,\delta\Phi_A(x),
\end{equation*}
其中
\begin{equation*}
 \mathfrak E_A
 \equiv
 \frac{\partial\Lag}{\partial\Phi_A}
 -\partial_\mu\frac{\partial\Lag}{\partial(\partial_\mu\Phi_A)}.
\end{equation*}
区域内部的 $\delta\Phi_A(x)$ 可以独立任取；要使 $\delta S=0$ 对所有变分都成立，只能有 $\mathfrak E_A(x)=0$。移项即得到\eqnrefs{eq:field-euler-lagrange-dp61}。
\end{derivation}

\begin{derivation}[从\eqnrefs{eq:noether-total-derivative-dp61}到\eqnrefs{eq:noether-current-general-dp61}]
链式法则给

\begin{align*}
 \delta\Lag
 &=\frac{\partial\Lag}{\partial\Phi_A}\delta\Phi_A
 +\frac{\partial\Lag}{\partial(\partial_\mu\Phi_A)}
 \partial_\mu\delta\Phi_A.
\end{align*}
对第二项做同样的乘积拆分，
\begin{align*}
 \delta\Lag
 &=\left[
 \frac{\partial\Lag}{\partial\Phi_A}
 -\partial_\mu\frac{\partial\Lag}{\partial(\partial_\mu\Phi_A)}
 \right]\delta\Phi_A\\
 &\quad+\partial_\mu\!\left[
 \frac{\partial\Lag}{\partial(\partial_\mu\Phi_A)}
 \delta\Phi_A\right].
\end{align*}
在满足 Euler--Lagrange 方程的经典解上，第一行的方括号为零。再代入 $\delta\Phi_A=\epsilon\Delta_A$：
\begin{equation*}
 \delta\Lag
 =\epsilon\,\partial_\mu\!\left[
 \sum_A\frac{\partial\Lag}{\partial(\partial_\mu\Phi_A)}
 \Delta_A\right].
\end{equation*}
另一方面，对称性假设给 $\delta\Lag=\epsilon\partial_\mu K^\mu$。两式相减，并利用常数 $\epsilon$ 可约去，得到
\begin{equation*}
 \partial_\mu\!\left[
 \sum_A\frac{\partial\Lag}{\partial(\partial_\mu\Phi_A)}
 \Delta_A-K^\mu\right]=0,
\end{equation*}
即\eqnrefs{eq:noether-current-general-dp61}。
\end{derivation}



